% !TeX root = 00Songbook02.tex
\begin{SBSong}{Stille Nacht, heilige Nacht}
    {$Id: 72StilleNacht.tex 689 2011-10-25 20:17:50Z LudwigVanBoll@gmail.com $}
    {First performed in the Nikolaus-Kirche, Oberndorf, Austria on 24 December 1818}
    {Joseph Mohr und Franz Xaver Gruber}
  \SBChord{G}
  \SBChord{D}
  \SBChord{C}
  \SBChord{D7}
  \SBDirection{Waltz timing  $\Downarrow \downarrow \downarrow $}
  \begin{SBVerse}
	\Ch{G}{Stille} Nacht, heilige Nacht,
	
	\Ch{D}{Alles} schl\"aft, \Ch{G}{einsam} wacht
	
	\Ch{C}{Nur} das traute \Ch{G}{hochheilige} Paar.
	
	\Ch{C}{Holder} Knabe im \Ch{G}{lockigen} Haar,

	\Ch{D7}{Schlaf} in himmlischer \Ch{G}{Ruh!}

    Schlaf in \Ch{D7}{himmlischer} \Ch{G}{Ruh!}
  \end{SBVerse}
  \begin{SBVerse}
	\Ch{G}{Stille} Nacht, heilige Nacht,
	
	\Ch{D}{Hirten} erst \Ch{G}{kund}gemacht

	\Ch{C}{Durch} der Engel \Ch{G}{Halleluja},

	\Ch{C}{T\"ont} es laut von \Ch{G}{fern} und nah:

	\Ch{D7}{Christ}, der Retter ist \Ch{G}{da!}

    Christ, der \Ch{D7}{Retter} ist \Ch{G}{da!}
  \end{SBVerse}
  \begin{SBVerse}
	\Ch{G}{Stille} Nacht, heilige Nacht,
	
	\Ch{D}{Gottes} Sohn, \Ch{G}{o} wie lacht

	\Ch{C}{Lieb'} aus deinem \Ch{G}{g\"ottlichen} Mund,

	\Ch{C}{Da} uns schl\"agt die \Ch{G}{rettende} Stund'.

	\Ch{D7}{Christ}, in deiner Ge\Ch{G}{burt!}

    Christ, \Ch{D7}{in} deiner Ge\Ch{G}{burt!}
  \end{SBVerse}
  \SBEnd{\Ch{D7}{Christ}, in deiner Ge\Ch{G}{burt!}
  
   Christ, in \Ch{D7}{deiner} Ge\Ch{G}{burt!}}
\end{SBSong}
